1994: \citet{hey_portability_1994}

first portability meeting


\newpage
1996: \citet{foster_performance_1996}

\begin{itemize}
    \item one of the first (?) PP papers
    \item performance and portability of three different Massively parallel processing computer systems (Intel Paragon, IBM SP2, Cray T3D)
    \item different algorithms used for the different architectures
    \item no single approach best on all problem sizes or hardware platforms
    \item performance portability enhanced by implementing multiple algorithms selectable at runtime
\end{itemize}


\newpage
2001: \citet{michalakes_performance-portability_2001}

\begin{itemize}
    \item performance portability on weather apps
    \item mainly also portability
    \item also weather code
    \item impact of storage order on performance and portability (emphasizes code portability)
    \item hauptsächlich TLB Analysen
\end{itemize}


\newpage
2007: \cite{zhu_performance_2007} (before the age of GPUs)

\begin{quotebox}{\citet{zhu_performance_2007}}
    If one parallel implementation of one application already achieves good performance on one platform, can this implementation also achieve the same or similar performance on other platforms?
\end{quotebox}

\begin{definition}[Performance Portablity $P^\mathcal{P}_n(b \rightarrow t$~\citet{zhu_performance_2007}]
Given a parallel program $\mathcal{P}$ developed on the \textbf{base system}, let $S^B_n$ denote the absolute speedup achieved by $\mathcal{P}$ on $n$ computing nodes of the \textbf{base system}, and $S^T_n$ denote the absolute speedup achieved by $\mathcal{P}$ on $n$ computing nodes of the \textbf{taregt system}, then the \textbf{performance portability} on $n$ computing nodes for porting the program $\mathcal{P}$ from the \textbf{base system} to the \textbf{target system} is:

$$P^\mathcal{P}_n(b \rightarrow t) = \frac{S^T_n}{S^B_n} \times \SI{100}{\percent}$$
\end{definition}

\begin{itemize}
    \item may be greater than $\SI{100}{\percent}$
    \item \enquote{Speedup, which is a measurement of the scalability of a parallel program, also reflects the performance portability of a parallel program in the sense that performance is portable with the increase of number of processors.}
\end{itemize}

\newpage
In 2014 \citeauthor{kunkel_performance_2014} defined performance portability as:
\begin{quotebox}{\citet{kunkel_performance_2014}}
    A code is performance portable if it can achieve a similar fraction of peak hardware performance on a range of different target architectures.
\end{quotebox}
As with the quote from \citeauthor{edwards_kokkos_2013} the problem is the objectivity of the definition namely how big the \enquote{\textit{similar fraction of}} should be. 
Although \citeauthor{kunkel_performance_2014} stated that for them it is \SI{20}{\percent} of the \enquote{best achievable performance} with respect to a hand-optimized native implementation, another person could argue that it should be \SI{15}{\percent} or \SI{25}{\percent} making this metric also highly subjective. 


\newpage
In 2013 in their original Kokkos paper \citeauthor{edwards_kokkos_2013} defined performance portability as:
\begin{quotebox}{\citet{edwards_kokkos_2013}}
    Our performance portability objective is to maximize the amount of user code that can be compiled for diverse devices and obtain the same (or nearly the same) performance as a variant of the code that is written specifically for that device.
\end{quotebox}
While this definition was already quite good, the problem is that \enquote{\textit{nearly the same}} can and will be interpreted differently by different people. 
Therefore, this definition is not feasible to use as an object performance portability metric.


\newpage 
In 2016 \citet{larkin_performance_2016} defined performance portability as:
\begin{quotebox}{\citet{larkin_performance_2016}}
    The same source code will run productively on a variety of different architectures.
\end{quotebox}
Again, \enquote{\textit{productively}} is highly subjective not only in a performance but also in a portable way. 
This definition can include that maintaining custom code paths for better performance may be ok as long as it is not too much work. 


\newpage
Also in 2016, the \gls{doe} conducted a meeting \cite{neely_doe_2016} specifically regarding performance portability. 
There, a few attendees were asked to define performance portability in their own words:
\begin{quotebox}{Rob Neely~\cite{neely_doe_2016}}
    For purposes of this meeting, it is the ability to run an application with acceptable performance across KNL and GPU-based systems with a single version of source code.
\end{quotebox}
\begin{quotebox}{Simon Pennycook~\cite{neely_doe_2016}}
    An application is performance portable if it achieves a consistent level of performance (e.g. defined by execution time or other figure of merit, not percentage of peak flops across platforms relative to the best known implementation on each platform.
\end{quotebox}
\begin{quotebox}{Adrian Pope, Vitali Morozov~\cite{neely_doe_2016}}
    Hard portability = no code changes and no tuning. Software portability = simple code mods with no algorithmic changes. Non-portable = algorithmic changes.
\end{quotebox}
\begin{quotebox}{David Richards~\cite{neely_doe_2016}}
    A code is performance portable when the application team says its performance portable!
\end{quotebox}

At the end of the \gls{doe} meeting, the attendees also created a working definition trying to incorporate the above definitions into one: 
\begin{quotebox}{\gls{doe} Working Definition~\cite{doe_meeting_attendees_definition_nodate}}
    An application is performance portable if it achieves a consistent ratio of the actual time to solution to either the best-known or the theoretical best time to solution on each platform with minimal platform specific code required.
\end{quotebox} % https://performanceportability.org/perfport/definition/

As with all other previous definitions, all of these quotes contain some objective part in one way or another.
However, to discuss performance portability objectively, these subjective definitions can not be used if it is expected that everyone should agree on the same values. 


\newpage
\cite{hoefler_scientific_2015}

\begin{itemize}
    \item experimental design (for reproducibility): application, input, compiler, runtime environment, machine, and measurement methodology
    \item interpretability: \enquote{We call an experiment interpretable if it provides enough information to allow scientists to understand the experiment, draw own conclusions, assess their certainty, and possibly generalize results}
\end{itemize}

Rules:
\begin{enumerate}
    \item When publishing parallel speedup, report if the base case is a single parallel process or best serial execution, as well as the absolute execution performance of the base case.
    \item Specify the reason for only reporting subsets of standard benchmarks or applications or not using all system resources.
    \item Use the arithmetic mean only for summarizing costs. Use the harmonic mean for summarizing rates.
    \item Avoid summarizing ratios; summarize the costs or rates that the ratios base on instead. Only if these are not available use the geometric mean for summarizing ratios.
    \item Report if the measurement values are deterministic. For nondeterministic data, report confidence intervals of the measurement.
    \item Do not assume normality of collected data (e.g., based on the number of samples) without diagnostic checking.
    \item Compare nondeterministic data in a statistically sound way, e.g., using non-overlapping confidence intervals or ANOVA.
    \item Carefully investigate if measures of central tendency such as mean or median are useful to report. Some problems, such as worst-case latency, may require other percentiles.
    \item Document all varying factors and their levels as well as the complete experimental setup (e.g., software, hardware, techniques) to facilitate reproducibility and provide interpretability.
    \item For parallel time measurements, report all measurement, (optional) synchronization, and summarization techniques.
    \item If possible, show upper performance bounds to facilitate interpretability of the measured results.
    \item Plot as much information as needed to interpret the experimental results. Only connect measurements by lines if they indicate trends and the interpolation is valid.
\end{enumerate}

Additional: 
\begin{itemize}
    \item report units unambiguously
    \item costs: metrics with a direct semantic (runtime, energy consumption, number of instructions)
    \item rates: derived from costs (\enquote{if the denominator has the primary semantic meaning, the harmonic mean provides correct results})
    \item ratios (unitless): should never be averaged
    \item statistical methods instead of arithmetic ones (samples must be normal distributed and either costs or rates): 
    \begin{itemize}
        \item standard variation
        \item coefficient of variation (good for performance consistency of a system over longer time)
        \item confidence interval of the mean
    \end{itemize}
    \item median and quartiles
    \item confidence intervals of the median
    \item recommend not to remove outliers
    \item analysis of variance (ANOVA)
    \item quantile regression
    \item factorial design
    \item warm-up, warm vs cold cache
    \item power-of-two best case for some algorithms (e.g., MPI\_Reduce)
    \item string vs weak scaling
    \item full benchmarks vs kernel benchmarks
    \item measure single events
    \item in general more than $5$ runs
    \item three cases of bounds models:
    \begin{itemize}
        \item ideal linear speedup
        \item serial overheads (Amdahls's law)
        \item parallel overheads
    \end{itemize}
    \item Points should only be connected if they indicate a trend and values between two points are expected to follow the line
\end{itemize}

% https://intel.github.io/p3-analysis-library/











\newpage
2016: \citet{pennycook_metric_2016}\\
2017 (preprint): \citet{pennycook_implications_2019}

Therefore, the first objective performance portability metric was introduced by \citeauthor{pennycook_implications_2019} at the end of 2016 published on arXiv~\cite{pennycook_metric_2016}, and finally submitted as a full paper at the beginning of 2017~\cite{pennycook_implications_2019}. 
It was the first formal definition that can be used to assess the performance portability of an application objectively.

\begin{definition}[Performance Portability \pp~(Pennycook et al.~\cite{pennycook_implications_2019})]\label{def:pp}
    Given a problem $\mathbf{p}$ which is solved by an application $\mathbf{a}$ and a set of platforms $\mathbf{H}$, the performance portability metric \pp is defined as:

    \[ \pp(a, p, H) = 
        \begin{cases}
            \frac{\abs{H}}{\sum\limits_{i \in H} \frac{1}{e_i(a, p)}} & \text{if }i\text{ is supported } \forall i \in H \\
            0                                                         & \text{otherwise}
        \end{cases} 
    \]

    where $\mathbf{e_i}$ is the performance efficiency of $\mathbf{a}$ solving $\mathbf{p}$ on platform $\mathbf{i \in H}$.
\end{definition}

Therefore, the performance portability metric \pp is the harmonic mean over the performance efficiency of an application solving a specific problem on all tested hardware platforms. 
The only difference to the harmonic mean is that \pp is also defined if the input data, i.e., the performance efficiency, includes a $0$. 
The \pp metric has the nice property that, per construction, it is in the interval $\interval{0}{1}$ where $0$ means the application did not run on all hardware platforms and $1$ means the best possible performance on all hardware platforms. 

Note that \pp heavily emphasizes the \textit{portability} aspect of performance portability: If only one hardware platform is not supported the result of \pp will directly be $0$ regardless of the performance on all other hardware platforms. 
Depending on your target audience, this property of \pp may be desirable or not. 


\newpage
\cite{pennycook_implications_2019}

\todo{irgendwie folgende Punkte noch unterbringen}
\begin{itemize}
    \item genaue Beschreibung von "platform": a specific hardware, OS, runtime system, etc. $\rightarrow$ hier \textbf{hardware}
    \item genaue Beschreibung von "application": software, die auf ihre Korrektheit überprüft werden kann $\rightarrow$ portability nicht nur definieren als "kann darauf laufen" sondern zusätzlich such "produziert korrektes resultat
    \item productivity schwierig zu messen/nicht objektiv, da sehr von skillset des jeweiligen Programmierers abhängig
    \item also formal Definition vor Pennycook rein?: Performance Portability on EARTH: A Case Study Across Several Parallel Architectures,
    \item arch: kann 100\% peak sein aber trotzdem nicht performant (z.B. bandwidth bound); z.B. via Roofline Model
    \item app: best known != best possible
    \item efficiency: 1 best possible; the lower, the worse
    \item 5 Kriterien an \pp aufnehmen?
    \item arch vs app : was rauslesen wenn eines klein, anderes groß
\end{itemize}

Performance Definition
\begin{quotebox}{\cite{pennycook_implications_2019}}
    Any measurable property of an application’s correct execution of a problem on a platform.
\end{quotebox}

Portability Definition
\begin{quotebox}{\cite{pennycook_implications_2019}}
    The ability of an application to execute a problem correctly on a given set of platforms.
\end{quotebox}

Performance Portability muss folgendes haben: 
\begin{itemize}
    \item beide obigen Definitionen nicht verletzen
    \item objektiv messbar, quantifizierbar und vergleichbar sein
\end{itemize}

Ihre Definition:
\begin{quotebox}{\cite{pennycook_implications_2019}}
    A measurement of an application’s performance efficiency for a given problem that can be executed correctly on all platforms in a given set.
\end{quotebox}


Current Problems in Performance Portability:
\begin{itemize}
    \item "abstraction and irreducible complexity": best Performance ohne spezialisierte Code-Pfade $\rightarrow$ hohe Abstraktionsebene, automatische Code-Generierung für performance
    \item "(still) no silver bullet": 
\end{itemize}

Approaches to Code Specialization
\begin{itemize}
    \item libraries
    \item DSL
    \item Portability and Productivity Frameworks
    \item Autotuning
    \item Application-specific abstractions
    \item performance portable algorithms

    \item we use frameworks + autotuning in the remainder of this work
    \item for more information on vor- und nachteile see paper \cite{pennycook_implications_2019}
\end{itemize}

Criteria: (verbatim quote)
\begin{enumerate}
    \item Be measured specific to a set of platforms of interest $\mathcal{H}$.
    \item Be independent of the absolute performance across $\mathcal{H}$.
    \item Be zero if a platform in $\mathcal{H}$ is unsupported, and approach zero as the performance of platforms in $\mathcal{H}$ approaches zero. 
    \item Increase if performance increases on any platform in $\mathcal{H}$.
    \item Be directly proportional to the sum of scores across $\mathcal{H}$.
\end{enumerate}


\newpage
Both formal performance portability definitions---\pp and \mpp---use a \textit{performance efficiency} $\mathbf{e_i(a, p)}$ in their definition which describes the relation between the achieved performance with respect to some proven expected performance. 
In general, two different definitions defining the performance efficiency of an application $\mathbf{a}$ solving problem $\mathbf{p}$ on the hardware $\mathbf{i \in H}$ are used. 

\begin{definition}[Architectural Efficiency]\label{def:arch_eff}
    The achieved performance as a fraction of the \enquote{peak} theoretical hardware performance, being it the achieved \si{\flop\per\second} or memory bandwidth.
\end{definition}

The \textit{architectural efficiency}

\begin{definition}[Application Efficiency~(\citet{satish_can_2012})]\label{def:app_eff}
    The achieved performance as a fraction of the \enquote{best} observed performance. 
\end{definition}

The \textit{application efficiency}











